\begin{exo}{}
	Pour chaque affirmation, indiquer si elle est \textbf{Vraie} ou \textbf{Fausse} en justifiant la réponse. (Si une affirmation est fausse, on donnera un contre-exemple).

	\begin{enumerate}[label=\arabic*), itemsep=5pt]
		\item Si un nombre appartient à l'ensemble $\mathbb{N}$, alors il appartient à l'ensemble $\mathbb{R}$.
		\item Si un nombre appartient à l'ensemble $\mathbb{D}$, alors il appartient à l'ensemble $\mathbb{Z}$.
		\item Tout nombre rationnel est un nombre réel.
		\item Si $x \in \mathbb{R}$ et $x \notin \mathbb{Q}$, alors $x$ est forcément un nombre décimal.
		\item Si un nombre entier relatif n'est pas un nombre naturel, alors il est strictement négatif.
		\item Le quotient de deux entiers relatifs est toujours un nombre décimal.
	\end{enumerate}
\end{exo}
